\subsection{Модуль lab3}

\subsubsection{MaxFlow\textless T\textgreater: алгоритм Форда--Фалкерсона (Эдмондс--Карп)}

\textbf{Вход:} \texttt{Graph<T> const\&}, исток \texttt{int s},
сток \texttt{int t}.

\textbf{Выход:} \texttt{T}~-- значение максимального потока.

\textbf{Описание:} Конструктор строит остаточный граф \texttt{cap}
из весовой матрицы; отрицательные веса берутся по модулю, потому что
пропускные способности должны быть неотрицательными.

Метод \texttt{bfs(s, t)} ищет кратчайший увеличивающий путь в остаточной
сети~-- ребро $(v, u)$ берётся только при $\texttt{cap}[v][u] > 0$.
Главный цикл \texttt{compute(s, t)}:
\begin{enumerate}
    \item ищет увеличивающий путь через BFS;
    \item находит узкое место~-- минимальную пропускную способность на пути;
    \item обновляет остаточный граф: прямые рёбра уменьшает, обратные
          увеличивает;
    \item прибавляет найденный поток к итогу.
\end{enumerate}

\begin{lstlisting}[style=cstyle, caption={Алгоритм Эдмондса--Карпа (maxflow.cpp)}]
template <typename T>
T MaxFlow<T>::compute(int s, int t) {
    T maxFlow = T{};
    while (true) {
        vector<int> path = bfs(s, t);
        if (path.empty()) break;

        T push = numeric_limits<T>::max();
        for (int i = 0; i + 1 < (int)path.size(); i++)
            push = min(push, cap[path[i]][path[i+1]]);

        for (int i = 0; i + 1 < (int)path.size(); i++) {
            cap[path[i]][path[i+1]] -= push;
            cap[path[i+1]][path[i]] += push;
        }
        maxFlow += push;
    }
    return maxFlow;
}
\end{lstlisting}

\textbf{Сложность:} $O(V \cdot E^2)$ по времени; $O(V^2)$ по памяти.

\subsubsection{MinCostFlow\textless T\textgreater: поток минимальной стоимости}

\textbf{Вход:} \texttt{Graph<T>}, матрица стоимостей
\texttt{vector<vector<T>>}, исток \texttt{s}, сток \texttt{t},
целевой поток \texttt{targetFlow}.

\textbf{Выход:} пара $(\text{достигнутый поток},\; \text{минимальная стоимость})$.

\textbf{Описание:} Конструктор заполняет остаточные пропускные способности
\texttt{cap} и матрицу стоимостей \texttt{cost}; обратные рёбра получают
противоположный знак стоимости.

На каждой итерации метод \texttt{findMinCostPath(s, t)}:
\begin{enumerate}
    \item строит вспомогательный граф \texttt{residual} из рёбер
          с ненулевой остаточной пропускной способностью;
    \item запускает \texttt{FloydWarshall<T>} для поиска пути
          минимальной стоимости от $s$ до $t$;
    \item возвращает путь и его стоимость.
\end{enumerate}

\begin{lstlisting}[style=cstyle, caption={Поиск пути минимальной стоимости (mincost.cpp)}]
template <typename T>
pair<vector<int>, T> MinCostFlow<T>::findMinCostPath(
    int s, int t) const
{
    Graph<T> residual(n);
    for (int i = 0; i < n; i++)
        for (int j = 0; j < n; j++)
            if (cap[i][j] > T{})
                residual.addEdge(i, j, cost[i][j]);

    FloydWarshall<T> fw(residual);
    fw.compute();
    const auto& dist = fw.getDistMatrix();
    const T INF_VAL = numeric_limits<T>::max() / 2;
    if (dist[s][t] >= INF_VAL) return {{}, T{}};
    return {fw.getPath(s, t), dist[s][t]};
}
\end{lstlisting}

По найденному пути определяется узкое место, поток направляется по пути,
остаточный граф обновляется. Итерации продолжаются до достижения
целевого потока $\lfloor 2/3 \cdot F_{\max} \rfloor$ или пока путь есть.

\textbf{Сложность:} $O(F \cdot V^3)$ по времени, где $F$~-- целевой поток
(каждая итерация запускает Флойда-Уоршалла); $O(V^2)$ по памяти.
